% ============================================================================
%  01_introduccion.tex — Introducción y punto de partida.
% ============================================================================
\section{Introducción}
\label{sec:intro}

La Actividad~3 cierra el ciclo de desarrollo del producto inteligente de
\emph{ACME\,S.A.}: sobre el sistema embebido construido en las dos actividades
previas se incorpora una funcionalidad de \textbf{instrumentación avanzada}
que perfecciona su eficiencia, fiabilidad y trazabilidad, alineándola con los
objetivos de digitalización industrial que persigue la empresa.

El \textbf{punto de partida} es una base limpia que fusiona en un único
\emph{sketch} sobre ESP32 los dos proyectos anteriores: el \textbf{clasificador
de cajas} de la Actividad~1 (HC-SR04, servo-compuerta, indicadores) y el
\textbf{ascensor inteligente de cinco plantas} de la Actividad~2 (mando IR
NEC con la librería \emph{Arduino-IRremote}~\cite{irremote_gh}, botones, PIR,
lazo ON-OFF de clima). La integración va más allá de la mera
convivencia en el mismo \codein{loop()}: cuando el clasificador confirma una
caja \emph{alta}, su altura se mapea a una planta del ascensor y se invoca
\codein{elevator\_request()} automáticamente, de modo que la caja desviada al
carril de ascensor en Act1 \emph{sube físicamente} por el ascensor de Act2. La
arquitectura modular de ambos proyectos --- pines disjuntos y APIs claras ---
permitió esa fusión con cambios mínimos y sin tocar el cableado de los
sensores compartidos.

Entre las tres líneas que propone el enunciado para perfeccionar el sistema
---control remoto, autodiagnóstico y supervisión inteligente--- se ha
seleccionado el \textbf{autodiagnóstico}. Es la opción que mejor encaja con
el reparto de evaluación (el Criterio~2, que vale el \SI{40}{\percent}, premia
la \emph{ampliación del circuito E/S}) y la que aporta más valor sobre la base
ya construida: introduce \textbf{sensores redundantes}, \textbf{detección de
averías por límites de rango}, \textbf{promediado de medidas} y
\textbf{vigilancia del consumo de actuadores}, todo ello con una
\textbf{reacción escalonada} que reproduce el comportamiento de un sistema de
instrumentación industrial moderno (avisar cuando hay reserva, dar alarma
cuando el fallo es crítico). La nueva capa se implementa en un único módulo
---\codein{diagnostics}--- que se interpone entre los sensores y el lazo de
control sin alterar los módulos heredados, cumpliendo así el principio de
\emph{mínima programación} (Criterio~3).

El resto de la memoria describe el hardware fusionado (§\ref{sec:hardware}),
la arquitectura del firmware (§\ref{sec:firmware}), el núcleo de la nueva
capa de instrumentación (§\ref{sec:autodiag}), la interfaz HMI y el logger
ampliado (§\ref{sec:hmi}), las pruebas reales en Wokwi (§\ref{sec:resultados})
y las conclusiones (§\ref{sec:conclusiones}).
